\statement{E.50}{Proved}{LM4 D.28 holds for the full pure three-cusp group}
Let \(p\) be a prime and \(\rho:C_p\to\mathcal P\) a
continuous homomorphism whose multiplier has
\(p\)-component one. Let \(\alpha\) be the unique normalized
actual representative of the image of a named generator,
and set
\[
 E=F\rtimes_\alpha C_p,\qquad
 S=\{1\}\times C_p,\qquad C=C_F(S).
\]
There is one \(F\)-conjugacy class of sections of
\(E\to C_p\). Each of the three peripheral normalizers
contains a section, and inclusion induces
\[
                     C^{(p)}\xrightarrow{\sim}F^{(p)}.
\]
Here \(C^{(p)}\) is the maximal pro-\(p\) quotient of
the actual full fixed subgroup. The statement is valid
for \(\mathcal P\), not merely for \(\widehat{GT}_0\).

\begin{proof}
E.51 proves the normalized homomorphic section without
using LM4 D.32. Hence \(\alpha^p=1\), and E.51's
Magnus argument makes \(\alpha\) identity on \(F^{(p)}\).
There is consequently a continuous map
\[
 E\longrightarrow F^{(p)}\times C_p,
 \qquad (g,t^j)\longmapsto(\bar g,t^j).
\]
Write \(c_1=x,c_2=y,c_3=(xy)^{-1}\) and
\(L_i=\overline{\langle c_i\rangle}\). By E.52,
\(N_F(L_i)=L_i\). Peripheral preservation thus gives
\(1\to L_i\to N_E(L_i)\to C_p\to1\).
Choose \(a_i\) over the named generator and write
\[
 a_i c_i a_i^{-1}=c_i^\lambda,
 \qquad a_i^p=c_i^{m_i},
 \qquad (\lambda-1)m_i=0.
\]
In the displayed pro-\(p\) quotient,
\(a_i\) has first coordinate \(\bar c_i^{v_i}\),
by the pro-\(p\) instance of E.52. Hence
\((m_i)_p=pv_i\). For
\(N_\lambda=1+\lambda+\cdots+\lambda^{p-1}\),
the equation \(N_\lambda u_i=-m_i\) is now solvable
in \(\widehat{\mathbb Z}\): at a prime \(q\) with
\(\lambda_q\ne1\) both the norm and \((m_i)_q\)
are zero; at every remaining prime the norm is \(p\),
and the only nonunit division was supplied above.
Thus \(a_ic_i^{u_i}\) has order \(p\) and defines
a peripheral section \(S_i\).

Let \(e_p\in\widehat{\mathbb Z}\) have coordinate
one at \(p\) and zero elsewhere. The group \(S_i\)
centralizes \(c_i^{e_p}\), whose mod-\(p\)
abelian image is the nonzero cusp vector.
Consequently \(H^1(C_F(S_i),\mathbb F_p)\) has
dimension at least one. For completeness, the
cohomological input used in the refereed proof is
Kochloukova--Zalesskii, Theorem~2 and Section~1.1,
equations (1)--(6), in the opened author manuscript.
It applies to full profinite groups with no
\(C_p^2\) and an open subgroup of finite
\(p\)-cohomological dimension.\cite{A_KZ2007}
Here every finite subgroup of \(E\) injects into
\(C_p\), and \(\operatorname{cd}_pF=1\).
For an order-\(p\) subgroup \(T\),
\(N_E(T)=C_F(T)\times T\). Restriction, K\"unneth,
and the Hochschild--Serre bound therefore give
\[
 \sum_{[T]}
 \bigl(1+\dim_{\mathbb F_p}H^1(C_F(T),\mathbb F_p)\bigr)
       =\dim_{\mathbb F_p}H^n(E,\mathbb F_p)\le3
       \qquad(n>1).
\]
The sum is finite: the restriction map is a dense
embedding from a finite-dimensional space into the
product, so it is surjective and that product is finite.

Two distinct cusp-section classes would contribute
at least four. They are therefore all conjugate.
Conjugating the \(x\)- and \(y\)-sections to one
representative places two elements with independent
mod-\(p\) abelian cusp vectors in its fixed subgroup.
Its contribution is at least three; the bound leaves
exactly this class and gives fixed-subgroup first
cohomology dimension two. Conjugacy under \(E\)
equals conjugacy under \(F\), since \(E=FT\).

The image of \(C\) in \(F^{(p)}\) is closed and
surjects onto its Frattini quotient, hence is all
of \(F^{(p)}\). Closed subgroups of the projective
group \(F\) are projective. The maximal pro-\(p\)
quotient of \(C\) is thus free pro-\(p\), with
rank two by the cohomology calculation. The induced
surjection between free pro-\(p\) groups of rank
two is an isomorphism by the Hopf property.
No \(S_3\)-symmetry equation or pentagon was used.
\end{proof}

The original LM4 D.28 register row restricted its
domain to \(\widehat{GT}_0\); E.50 records the domain
actually needed in LM4 D.32. Once R1 is established,
the displayed hypothesis forces \(\rho\) to be
trivial, so this conditional lemma has no additional
nontrivial torsion examples. Its proof above is given
independently precisely to preserve the logical
dependency used to establish R1.

\statement{E.51}{Proved}{The normalized section and the independent Magnus lemma remove the citation cycle}
The pure peripheral outer group \(\mathcal P\)
has a unique normalized homomorphic section into
\(\operatorname{Aut}(F)\) of the form
\[
 x\longmapsto x^\lambda,
 \qquad y\longmapsto f^{-1}y^\lambda f,
 \qquad f\in F'.
\]
For every prime \(p\), a finite-order continuous
automorphism of a finitely generated free pro-\(p\)
group which is identity on integral abelianization
is identity. This lemma is independent of LM4
D.25's final conclusion and of LM4 D.32.

\begin{proof}
For the first assertion, a peripheral outer class
has common multiplier \(\lambda\), since the
abelian relation \([x]+[y]+[z]=0\) equates the
three exponents. After an inner adjustment its
representative sends \(x\) to \(x^\lambda\) and
\(y\) to \(k^{-1}y^\lambda k\). Write
\([k]=(a,b)\in\widehat{\mathbb Z}^2\).
Conjugating that representative by \(x^a\) and
replacing its \(y\)-conjugator on the left by
\(y^{-b}\) gives
\(f=y^{-b}kx^{-a}\in F'\).
If two normalized representatives are inner-equivalent,
the inner conjugator centralizes \(x^\lambda\),
so is \(x^v\) by E.52. Their \(y\)-conjugators
then differ by left multiplication by \(y^w\)
and right multiplication by \(x^{-v}\).
The two zero abelian coordinates force \(v=w=0\).
Composition of normalized representatives has
conjugator \(f\alpha(g)\in F'\), proving the
homomorphism assertion as well as uniqueness.

For the second assertion, let the free basis be
\(x_1,\ldots,x_d\). Its integral Magnus map is
the continuous injective homomorphism into
\[
 \mathbb Z_p\langle\!\langle X_1,\ldots,X_d
               \rangle\!\rangle^{\times},
 \qquad x_i\longmapsto1+X_i.
\]
Continuity uses the \((p,X_1,\ldots,X_d)\)-adic
topology; injectivity follows already after reduction
modulo \(p\) from the faithful free pro-\(p\)
Magnus map. An automorphism identity on integral
abelianization substitutes
\(X_i\mapsto X_i+R_i\), with every \(R_i\)
of homogeneous degree at least two.
If some \(R_i\ne0\), choose the least occurring
degree \(m\). Induction on \(n\) shows that the
degree-\(m\) correction in the \(n\)-th iterate
is exactly \(n\) times the original correction:
any new substitution into a degree-\(m\) term
changes only larger degrees. The homogeneous
coefficient module is torsion-free over
\(\mathbb Z_p\), so no positive iterate can be
identity. This contradiction proves the lemma,
including \(p=2\).
\end{proof}

The acyclic dependency is now: E.52 supplies the
free-factor centralizers; E.51 supplies normalization
and the Magnus torsion lemma; E.50 supplies the
prime-order section calculation; the separate
three-cusp finite-cover argument in refereed
LM4 D.32 proves R1. Only after that may LM4 D.25
deduce triviality of its full multiplier-one
finite-order solutions. Its final R1-dependent
sentence must not be cited as an input to D.32.

\statement{E.52}{Known}{The precise Herfort--Ribes citation for the free peripheral reduction}
For \(F=\widehat F(x,y)\), each peripheral cyclic
group \(L=\overline{\langle u\rangle}\), where
\(u=x,y,(xy)^{-1}\), is malnormal, and
\[
 N_F(L)=L,
 \qquad C_F(u^a)=L
       \quad(0\ne a\in\widehat{\mathbb Z}).
\]
The analogous statements hold in \(F^{(p)}\),
with \(0\ne a\in\mathbb Z_p\). The precise
source in the free-factor reduction of LM4 A.7
is Herfort--Ribes (1985), \emph{Theorem B'},
Section~4, printed page~159. \cite{A_HR1985}

\begin{proof}
Theorem B' states that a factor of a free
pro-\(\mathcal C\) product intersects each of
its conjugates outside itself trivially, and
identifies the centralizer of a nonidentity
factor element with its centralizer in that
factor. Use the full formation of finite
groups, or the formation of finite \(p\)-groups.
Each of the three displayed \(u\)'s is a free
basis element: use \((x,y)\), \((y,x)\), or
\((x,(xy)^{-1})\), respectively. Hence
\(F=L\amalg L'\), and Theorem B' applies.
The power \(u^a\) is nonidentity because
\(L\simeq\widehat{\mathbb Z}\), or
\(\mathbb Z_p\). A normalizing element has
nontrivial factor intersection and therefore
belongs to \(L\). The reverse inclusions
follow from procyclicity.
\end{proof}

This citation correction concerns the free-factor
argument, which covers the pants group used in
LM4 D.28 and D.32. For a once-punctured surface
the boundary word is generally not primitive;
the separate peripheral commensurator theorem
used in LM4 A.7 cannot be replaced by this
free-factor citation. The locations in the
manuscript remain \S13 bis, folio~57 verso,
scan~118, and \S35, folios~253--254,
scans~499--501; E.50 also concerns \S31,
folios~198--199 verso, scans~391--394.
